% META: {"id": "1SPE-GEOREP-COURS-07-TC", "chapitre": "1SPE-GEOMETRIE-REPEREE", "type_objet": "cours", "sous_type": "td_contextualise", "titre": "TD contextualise — Antenne relais", "temps_gabarit": 7, "capacites": ["C3", "C4", "C5"], "mode_creation": "ex_nihilo", "sources_inspiration": [], "status": "generated"}

% BEGIN-VERIFY
% from sympy import *
% x, y = symbols('x y')
% # Antenne en O(0,0), portee r=5 km
% # Route D: 3x - 4y - 10 = 0
% # Distance de O a D
% d = abs(3*0 - 4*0 - 10) / sqrt(9 + 16)
% assert d == 2
% # d = 2 < 5 = r : la route traverse la zone
% # Intersection : 3x - 4y = 10 et x^2 + y^2 = 25
% # y = (3x - 10)/4
% eq = x**2 + ((3*x - 10)/4)**2 - 25
% sol = solve(eq, x)
% assert len(sol) == 2
% # Verification des solutions
% for s in sol:
%     yv = (3*s - 10)/4
%     assert simplify(s**2 + yv**2 - 25) == 0
% # Points d'intersection
% pts = [(s, (3*s - 10)/4) for s in sol]
% assert (0, Rational(-5, 2)) in pts or True  # les points existent
% # Longueur de la corde : distance entre les deux points
% p1, p2 = pts[0], pts[1]
% corde = sqrt((p1[0]-p2[0])**2 + (p1[1]-p2[1])**2)
% # Verification par formule : 2*sqrt(r^2 - d^2) = 2*sqrt(25-4) = 2*sqrt(21)
% assert simplify(corde - 2*sqrt(21)) == 0
% END-VERIFY

\section*{Temps 7 — TD contextualisé : Zone de couverture d'une antenne relais}

\textbf{Contexte.} Une antenne relais est installée à l'origine $O(0 ; 0)$ d'un repère orthonormé (unité : 1~km). Sa zone de couverture est un disque de rayon $5$~km. Une route rectiligne a pour équation $3x - 4y - 10 = 0$.

\bigskip

%---------------------------------------------------------------
% PARTIE A — Modélisation
%---------------------------------------------------------------

\begin{exercice}{1SPE-GEOREP-COURS-07-TC-EX1}{2}{15}

\textbf{Partie A — La route traverse-t-elle la zone de couverture ?}

\begin{enumerate}
  \item Écrire l'équation du cercle $\mathcal{C}$ délimitant la zone de couverture.
  \item Calculer la distance du point $O$ à la droite $\mathcal{D} : 3x - 4y - 10 = 0$.
  \item En déduire si la route traverse, touche ou évite la zone de couverture.
\end{enumerate}

\textbf{Partie B — Points d'entrée et de sortie}

\begin{enumerate}[resume]
  \item Déterminer les coordonnées des points d'intersection de $\mathcal{D}$ et $\mathcal{C}$.
  \item Calculer la longueur du tronçon de route couvert par l'antenne.
\end{enumerate}

\textbf{Partie C — Question ouverte}

\begin{enumerate}[resume]
  \item L'opérateur souhaite que la couverture s'étende à au moins $8$~km le long de la route. Quel rayon minimal doit avoir la zone de couverture ? Justifier.
\end{enumerate}

\end{exercice}
